\title{FlashAttention-3: Fast and Accurate Attention with Asynchrony and Low-precision}

\begin{abstract}
The attention mechanism constitutes the central computational element in the Transformer architecture, and remains a limiting factor for scaling large language models and applications requiring extended contexts. \textsc{FlashAttention} introduced optimizations for accelerating attention calculations on GPUs by reducing the frequency of memory accesses, but prior versions did not fully exploit features available in recent hardware releases. In particular, \textsc{FlashAttention-2} utilized only 35\% of the H100 GPU's capacity. This work presents three strategies to enhance attention performance on Hopper GPUs: (1) leveraging the asynchronous nature of Tensor Cores and TMA through warp-specialization to simultaneously compute and transfer data, (2) alternating block-level matrix multiplication and softmax operations, and (3) employing block quantization alongside incoherent processing methods that take advantage of native FP8 support for low-precision calculations. The resulting approach, named \textsc{FlashAttention-3}, achieves a 1.5–2.0$\times$ acceleration on H100 GPUs, with FP16 computations reaching up to 740 TFLOPs/s (corresponding to 75\% hardware utilization) and FP8 calculations attaining nearly 1.2 PFLOPs/s. Additionally, evaluation shows that FP8 \textsc{FlashAttention-3} yields numerical errors 2.6$\times$ lower than those found in standard FP8 attention implementations.
\end{abstract}

\section{Introduction}
\label{sec:intro}

Within the Transformer architecture~\citep{vaswani2017attention}, the primary computational constraint arises from the attention mechanism, specifically due to the quadratic cost associated with calculating self-attention scores across queries and keys relative to the length of the input sequence. Enhancing the scalability of attention mechanisms to handle more extensive contexts has the potential to enable novel functionalities such as cross-document modeling and reasoning~\citep{guo2021longt5,shaham2022scrolls,peng2023yarn}, comprehensive codebase analysis~\citep{roziere2023code, li2023starcoder}, support for additional modalities—including high-resolution images~\citep{chen2022scaling}, audio~\citep{gulati2020conformer}, and video~\citep{ho2022video}—as well as emerging applications like interaction histories~\citep{sun2019bert4rec} and extended agent planning~\citep{yao2022react}. This challenge has catalyzed substantial research interest in accelerating attention operations for long-context scenarios through methods such as approximating attention computations~\citep{katharopoulos2020transformers,choromanski2020rethinking, tay2020efficient}, improving software efficiency~\citep{rabe2021self,dao2022flashattention,kwon2023efficient}, and exploring alternative architectural designs~\citep{peng2023rwkv,sun2023retentive,gu2023mamba}.

This study extends the contributions of \citet{dao2022flashattention}, who devised precise attention algorithms by incorporating insights from the GPU execution model and hardware-specific properties into their algorithmic frameworks. Dao et al. presented \textsc{FlashAttention}, which uses a unique tiling approach to enable efficient attention computation by consolidating all necessary operations into one GPU kernel, thereby avoiding costly intermediate accesses to global memory.

Building upon this foundation, \citet{dao2023flashattention2} further adapted the algorithm as \textsc{FlashAttention-2}, allowing for parallelization along the sequence length axis. This redesign executes the core portion of the forward pass over tiled sections of the key and value matrices, enhancing both workload distribution and GPU occupancy. Nonetheless, our analysis shows that \textsc{FlashAttention-2} continues to lag behind highly-tuned matrix multiplication (GEMM) kernels on newer hardware, such as achieving only about 35\% efficiency compared to 80-90\% on the Hopper H100 GPU. This gap can be partly explained by implementation factors, notably the lack of Hopper-specific instructions replacing those tailored for Ampere when leveraging Tensor Cores. Previous research, including ThunkerKitten~\citep{spector2024thunder} and cuDNN 9~\citep{cudnn9}, demonstrates that utilizing Hopper-specific operations and tile-oriented abstractions facilitates faster attention calculations and streamlines the codebase.

At a foundational level, the approach utilized by \textsc{FlashAttention-2} relies on a basic synchronous paradigm, deliberately omitting considerations for asynchrony and reduced numerical precision within its algorithmic framework. Asynchrony emerges from advanced hardware designs intended to optimize critical machine learning operations—such as dedicated hardware components for matrix multiplications (Tensor Cores) and memory access (Tensor Memory Accelerator -- TMA), which operate independently from general CUDA cores tasked with handling logical, integer, and floating-point processes. The utilization of lower precisions (like FP8 on Hopper and FP4 on Blackwell, following the FP16 standard in Pascal (2017) and BF16 in Ampere (2020)) represents an established method for achieving significantly increased computational throughput at constant energy and silicon footprint. A discussion of Hopper’s features relevant to these advances can be found in \cref{subsec:hardware}. The primary challenge involves reengineering \textsc{FlashAttention-2} so that it can fully exploit these specialized hardware functionalities: implementing asynchrony necessitates enabling concurrent operations between matmul and softmax—even though dependencies exist between their outputs—while low-precision arithmetic mandates careful strategies to suppress quantization artifacts, particularly in light of atypical values found within LLMs~\citep{dettmers2208llm, sun2024massive}.

With this objective in mind, we introduce \textsc{FlashAttention-3}, which integrates three novel concepts designed to enhance performance on the latest GPU architectures:\footnote{Our findings are discussed with reference to NVIDIA's Hopper architecture, though the algorithm remains applicable to any GPU architecture offering strong support for asynchronous execution and low-precision operations.}

\begin{enumerate}

\item \textbf{Producer-Consumer asynchrony:} We introduce a warp-specialized software pipelining strategy that leverages asynchronous data transfer and Tensor Core computation by dividing the roles of data producers and consumers across distinct warps. This approach broadens the algorithm’s capability to mask memory access and instruction issuing delays.

\item \textbf{Hiding softmax under asynchronous block-wise GEMMs:} To improve performance, we schedule the relatively slower non-GEMM operations associated with softmax—such as floating-point multiply-add and exponentiation—to coincide with asynchronous WGMMA-based GEMM computations. This entails refactoring the \textsc{FlashAttention-2} algorithm to sidestep certain serial dependencies between softmax computations and GEMM calls. For instance, in the two-stage variant, while softmax processes a block of the scores matrix, WGMMA proceeds asynchronously to compute the next block.

\item \textbf{Hardware-accelerated low-precision GEMM:} We modify the forward pass to make use of FP8 Tensor Cores during GEMM, achieving nearly a twofold increase in observed TFLOPs/s. Accommodating the distinct layout requirements of WGMMA—particularly regarding the placement of FP32 accumulator and FP8 operand matrices in memory—is necessary. To address the precision reduction associated with FP8, we employ block quantization and incoherent computation strategies.

To empirically assess the effectiveness of our approach, we conducted a comprehensive benchmarking of \textsc{FlashAttention-3} on the H100 SXM5 GPU using a variety of configurations. Our evaluation demonstrates: (1) using FP16, \textsc{FlashAttention-3} achieves a forward pass acceleration of 1.5-2.0$\times$ compared to \textsc{FlashAttention-2} (with throughput peaking at 740 TFLOPs/s) and a backward pass acceleration of 1.5-1.75$\times$; (2) FP8 yields performance near 1.2 PFLOPs/s; and (3) for extended sequence lengths, FP16 attains superior performance while FP8 remains competitive\footnote{Specifically, for head dimension 64, \textsc{FlashAttention-3} FP8 surpasses alternatives, whereas with head dimensions 128 and 256, FP8 is on par in the absence of causal masking but trails with causal masking.} with NVIDIA’s cuDNN attention implementation. Furthermore, we confirm that FP16 \textsc{FlashAttention-3} produces equivalent numerical errors as \textsc{FlashAttention-2}, and outperforms conventional attention operations since intermediate computations (such as softmax normalization) utilize FP32 precision. Lastly, FP8 \textsc{FlashAttention-3}, when paired with block quantization and incoherent computation, delivers 2.6$\times$ greater accuracy compared to typical attention employing per-tensor quantization in the presence of outlier features.

We have made \textsc{FlashAttention-3} openly available under a permissive license\footnote{\textsc{FlashAttention-3}
  is available at \texttt{https://github.com/Dao-AILab/flash-attention}}, and intend to incorporate it into the PyTorch and Hugging Face ecosystems to maximize its accessibility for both researchers and developers.

\section{Background: Multi-Head Attention and GPU Characteristics}
\label{sec:background}

\subsection{Multi-Head Attention}
\label{subsec:multi_head_attn}

Let $\mathbf{Q}, \mathbf{K}, \mathbf{V} \in \mathbb{R}^{N \times d}$ be the query, key and value input sequences associated to a single head, where $N$ is the sequence length and $d$ is the head dimension. Then the attention output $\mathbf{O}$ is computed as:

\begin{equation*}
  \mathbf{S} = \alpha \mathbf{Q} \mathbf{K}^\top \in \mathbb{R}^{N \times N}, \quad \mathbf{P} = \mathrm{softmax}(\mathbf{S}) \in \mathbb{R}^{N \times N}, \quad \mathbf{O} = \mathbf{P}\mathbf{V} \in \mathbb{R}^{N \times d},
\end{equation*}

Here, $\mathrm{softmax}$ operates along each row, and it is common to choose the scaling parameter $\alpha = 1/\sqrt{d}$. To address potential numerical issues with exponentiation, we adjust $\mathbf{S}$ by removing $\mathrm{rowmax}(\mathbf{S})$. In the case of multi-head attention (MHA), independent projections for queries, keys, and values are assigned to each head, allowing this operation to proceed concurrently across all heads and batch instances, which yields the complete output tensor.

Now let $\phi$ be a scalar loss function and let $\mathbf{d}(-) = \partial \phi / \partial (-)$ be notation for the gradient. Given the output gradient $\mathbf{dO} \in \mathbb{R}^{N \times d}$, we compute $\mathbf{dQ}$, $\mathbf{dK}$, and $\mathbf{dV}$ according to the chain rule as follows:

\begin{align*}
  \mathbf{dV} &= \mathbf{P}^\top \mathbf{dO} \in \mathbb{R}^{N \times d} \\
  \mathbf{dP} &= \mathbf{dO} \mathbf{V}^\top \in \mathbb{R}^{N \times N} \\
  \mathbf{dS} &= \mathrm{dsoftmax} (\mathbf{dP}) \in \mathbb{R}^{N \times N} \\
  \mathbf{dQ} &= \alpha \mathbf{dS} \mathbf{K} \in \mathbb{R}^{N \times d} \\
  \mathbf{dK} &= \alpha \mathbf{dS}^\top \mathbf{Q} \in \mathbb{R}^{N \times d},
\end{align*}

In this context, we observe that $\mathbf{d}s = (\mathrm{diag}(p) - p p^\top)\mathbf{d}p$, where $p$ is the result of applying $\mathrm{softmax}$ to the vector $s$. We use $\mathrm{dsoftmax}(\mathbf{dP})$ to denote the application of this operation across each row. Moreover, the backward pass in MHA leverages this computation, which is efficiently executed in parallel over multiple heads and batches.

\subsection{GPU hardware characteristics and execution model}
\label{subsec:hardware}

We outline the elements of the GPU execution model pertinent to \textsc{FlashAttention-3}, emphasizing the NVIDIA Hopper architecture as a specific example embodying this model.

\paragraph{Memory hierarchy:} The GPU implements a layered memory architecture, where data storage units are arranged such that smaller, faster memories possess higher bandwidth and lower capacity (\cref{tab:gpu-hierarchy})\footnote{\citet{luo2024benchmarking} provides a measurement of 128 bytes per clock cycle per SM for shared memory bandwidth; to obtain the total, we multiply by 132 SMs and a boost clock of 1830 MHz.}. The primary form of external DRAM, often called global memory (GMEM) or HBM, can be accessed by all streaming multiprocessors (SMs). Any data loaded from GMEM is automatically stored in an intermediate on-chip L2 cache. Furthermore, every SM is equipped with a compact, programmer-controlled, multi-banked shared memory (SMEM). At the lowest level of the hierarchy is the SM’s register file.

\paragraph{Thread hierarchy:} In the GPU programming paradigm, threads are structured into several layers of abstraction. These layers range from the most granular, individual threads, to warps, which consist of 32 threads, followed by warpgroups that aggregate 4 sequential warps. Higher up, threadblocks—also referred to as cooperative thread arrays (CTAs)—serve as larger organizational units. On Hopper architectures, threadblock clusters further combine threadblocks, and the top level is the grid, encompassing all lower levels.

The two hierarchies exhibit a strong interconnection. Threads that belong to a single CTA are concurrently scheduled to run on the same SM, and similarly, CTAs within one cluster are executed together on the same GPC. All threads inside a CTA have direct access to SMEM, while each thread is limited to a maximum of 256 private registers (RMEM) exclusively accessible by itself.

\paragraph{Asynchrony and warp-specialization:}

GPUs function as throughput-oriented processors that depend on concurrent execution and asynchronous operations to mask delays related to memory access and computation. Hopper introduces the Tensor Memory Accelerator (TMA), a specialized hardware component, for asynchronous memory transfers between GMEM and SMEM \cite[\S7.29]{cuda}. Additionally, differing from earlier GPU architectures like Ampere, Hopper’s Tensor Core—accessible through the warpgroup-level WGMMA instruction \cite[\S9.7.14]{ptx}—also operates asynchronously and is capable of obtaining its input data directly from shared memory.

Enabling asynchrony through hardware permits the design of warp-specialized kernels, wherein individual warps within a CTA are assigned exclusive producer or consumer tasks, focusing solely on data transfer or computational operations. This separation typically enhances the compiler's capacity to create more effective instruction scheduling \citep{warp-specialization-2011}. Furthermore, Hopper introduces the capability for warpgroups to dynamically redistribute their register usage through \verb|setmaxnreg| \cite[\S9.7.17.1]{ptx}, enabling warps responsible for MMA operations to access a greater portion of RMEM, in contrast to those executing TMA—an activity that can be handled with just a single thread.

\paragraph{Low-precision number formats:}
\label{sec:low-precision-gpu}
Contemporary GPUs incorporate dedicated hardware for efficient processing of low-precision arithmetic. For instance, the WGMMA instruction enables exploitation of the FP8 Tensor Cores within Hopper architecture, providing double the throughput per SM relative to computations executed in FP16 or BF16 formats.

To successfully use FP8 WGMMA, it is crucial to be aware of the operand layout restrictions. Consider a GEMM operation computing $A \times B^{\top}$, where $A$ is an $M\times K$ matrix and $B$ an $N\times K$ matrix. We classify $A$ or $B$ as \emph{mn-major} if elements are contiguous along the $M$ or $N$ dimension, and as \emph{k-major} when contiguity is along the $K$ dimension. FP16 WGMMA is flexible: it allows both mn-major and k-major inputs from SMEM. Conversely, FP8 WGMMA restricts input format to only k-major layouts. This limitation becomes particularly challenging in cases like attention mechanisms, where there is a need to fuse consecutive GEMM operations within a single kernel. Here, the differing layouts between FP32 accumulators and FP8 operands can impede seamless chaining of dependent FP8 WGMMAs.

Within the scope of attention mechanisms, these constraints regarding layout require adjustments to the structure of an FP8 algorithm. The details of these changes are provided in \cref{sec:algofp8}.

\subsection{Standard Attention and Flash Attention} As established by \citet{dao2022flashattention}, \textbf{standard attention} refers to the conventional GPU-based implementation in which the intermediate matrices $\mathbf{S}$ and $\mathbf{P}$ are explicitly stored in HBM. \textsc{FlashAttention} introduced a strategy to circumvent the costly memory traffic associated with these intermediates by implementing a local variant of the softmax operation, thus allowing the entire attention mechanism to be executed within a single kernel. This local softmax mechanism is executed in lines \ref{code-ws:softmax_start}-\ref{code-ws:softmax_end} of the consumer mainloop presented in \cref{alg:flash3_wgmma_ws_only}, including the necessary rescaling steps for blocks of $\mathbf{O}$. A straightforward proof that this approach yields the correct computation of $\mathbf{O}$ is provided in \cite[\S 2.3.1]{dao2023flashattention2}.

\section{FlashAttention-3: Algorithm}
\label{sec:algo}

This section outlines the details of the \textsc{FlashAttention-3} algorithm. Our exposition centers on the forward pass, while the method for the backward pass is provided in~\cref{sec:algo_ws_bwd}. We start by describing the integration of warp-specialization and a circular SMEM buffer with the underlying framework established in \textsc{FlashAttention-2}. Next, we elaborate on leveraging the asynchrony intrinsic to WGMMA to construct a two-stage pipeline in which GEMM and softmax computations are overlapped. Lastly, we present the adjustments necessary for handling FP8, including compliance with layout requirements and enhancements to accuracy through block quantization and incoherent processing strategies.

\subsection{Producer-Consumer asynchrony through warp-specialization and pingpong scheduling}
\label{sec:algo_ws}

\paragraph{Warp-specialization}
Similar to \textsc{FlashAttention-2}, the forward pass in \textsc{FlashAttention-3} exhibits a high degree of parallelism across the batch size, head count, and the length of the query sequence. Therefore, analyzing the algorithm at the CTA level is adequate; here, the process applies to a query tile $\mathbf{Q}_i$ and generates the respective output tile $\mathbf{O}_i$. For clarity, we introduce the warp-specialization strategy employing a circular SMEM buffer, excluding for now any overlap with GEMM-softmax operations. Define $d$ as the head dimension and $N$ as the sequence length, and choose a query block size $B_r$ such that the query matrix $\mathbf{Q}$ is partitioned into $T_r = \lceil \frac{N}{B_r} \rceil$ blocks: $\mathbf{Q}_1, .., \mathbf{Q}_{T_r}$.

\begin{algorithm}[H]
    \caption{\small\label{alg:flash3_wgmma_ws_only}\textsc{FlashAttention-3} forward propagation \textbf{excluding} intra-consumer overlapping -- CTA layout}
    \begin{algorithmic}[1]
\REQUIRE Matrices $\mathbf{Q}_i \in \mathbb{R}^{B_r \times d}$ as well as $\mathbf{K}, \mathbf{V} \in \mathbb{R}^{N \times d}$ stored in HBM, with a key block size parameter $B_c$ and total key block count $T_c = \lceil \frac{N}{B_c} \rceil$.
\STATE Instantiate a pipeline control structure to handle synchronization using a circular SMEM buffer with $s$ stages.
\IF {executing in the producer warpgroup}
\STATE Release a specified number of registers.
\STATE Start a transfer of $\mathbf{Q}_i$ from HBM into shared memory.
\STATE Upon successful load, signal consumers regarding the readiness of $\mathbf{Q}_i$ in shared memory.
\FOR{$0 \le j < T_c$}
    \STATE Block until the consumer has finished with the $(j\,\%\,s)$-th stage of the buffer.
    \STATE Initiate transfers for $\mathbf{K}_j, \mathbf{V}_j$ from HBM into shared memory at the $(j\,\%\,s)$-th buffer stage.
    \STATE Once the load is finished, inform consumers about the new availability of $\mathbf{K}_j, \mathbf{V}_j$.
\ENDFOR
\ELSE 
\STATE Allocate registers as determined by the consumer warp count.
\STATE Locally, initialize $\mathbf{O}_i = (0) \in \mathbb{R}^{B_r \times d}$, and set $\ell_i, m_i = (0), (-\infty) \in \mathbb{R}^{B_r}$.
\STATE Stall until $\mathbf{Q}_i$ is loaded into shared memory.
\FOR{$0 \le j < T_c$}
\STATE Hold until $\mathbf{K}_j$ is ready in shared memory.
\STATE Perform $\mathbf{S}_i^{(j)} = \mathbf{Q}_i \mathbf{K}_j^T$ (using SS-GEMM). Commit and synchronize. \label{alg:ws_only_gemm1}
\STATE Store the prior value $m_i^{\mathrm{old}} = m_i$ and update $m_i = \mathrm{max}(m_i^{\mathrm{old}}, \mathrm{rowmax}(\mathbf{S}_i^{(j)}))$. \label{code-ws:softmax_start}
\STATE Calculate $\widetilde{\mathbf{P}}_i^{(j)} = \mathrm{exp}(\mathbf{S}_i^{(j)} - m_i)$; also update $\ell_i = \mathrm{exp}(m_i^{\mathrm{old}} - m_i) \ell_i + \mathrm{rowsum}(\widetilde{\mathbf{P}}_i^{(j)})$. \label{code-ws:softmax_end}
\STATE Pause until $\mathbf{V}_j$ becomes available in shared memory.
\STATE Update $\mathbf{O}_i = \mathrm{diag}(\mathrm{exp}(m_i^{\mathrm{old}} - m_i))^{-1} \mathbf{O}_i + \widetilde{\mathbf{P}}_i^{(j)} \mathbf{V}_j$ via RS-GEMM; commit and synchronize. \label{alg:ws_only_gemm2}
\STATE Mark the $(j\,\%\,s)$-th buffer stage as free for producer usage.
\ENDFOR
\STATE Finalize with $\mathbf{O}_i = \mathrm{diag}(\ell_i)^{-1} \mathbf{O}_i$ and $L_i = m_i + \log(\ell_i)$.
\STATE Write the computed $\mathbf{O}_i$ and $L_i$ back to HBM for the $i$-th block of $\mathbf{O}$ and $L$.
\ENDIF
\end{algorithmic}
\end{algorithm}

In our deployment of \cref{alg:flash3_wgmma_ws_only} on Hopper, (de)allocation tasks are handled using \verb|setmaxnreg|, while TMA is employed for loading $\mathbf{Q}_i$ along with $\{ \mathbf{K}_j, \mathbf{V}_j \}_{0 \leq j < T_c}$. The consumer mainloop relies on WGMMA to perform the GEMMs, with either the SS or RS prefix specifying whether the first operand resides in shared memory or in the register file. To understand the execution pattern of \cref{alg:flash3_wgmma_ws_only}, it is important to recognize that TMA loads operate asynchronously, meaning their issuance is independent from the completion of prior loads. Furthermore, in the producer mainloop, no wait instructions are executed during the initial $s$ iterations as the buffer is being populated.

\paragraph{Pingpong scheduling} Leveraging the asynchronous execution capabilities of WGMMA and TMA, as well as warp-specialization, enables simultaneous execution of the softmax calculations for one warpgroup with GEMM operations in another. This approach is particularly advantageous given the significant disparity in performance between matmul and non-matmul operations on contemporary hardware accelerators. For instance, the H100 SXM5 GPU boasts a matmul throughput of 989 TFLOPS at FP16 precision, in contrast to just 3.9 TFLOPS for special functions like the exponential\footnote{According to the CUDA programming guide, each streaming multiprocessor (SM) can execute 16 special function operations per clock cycle. By multiplying 16 by the number of SMs (132) and the clock speed (1830 MHz), we derive the 3.9 TFLOPS rate for these functions.}, which are essential for softmax. In FP16 attention forward passes with a head dimension of 128, matmul FLOPS outnumber exponential operations by a factor of 512; however, exponential functions operate at 256 times lower throughput. As a result, exponentials may occupy up to 50\% of the computational cycle that matmuls require. The disparity further increases for FP8 computations: matmul throughput is doubled, but the exponential throughput remains unchanged.

Because the exponential computation is handled by a distinct hardware component (the multi-function unit), it is optimal to overlap the execution of the exponential calculation with the matmul operations on the Tensor Cores. To achieve this concurrency, we employ synchronization barriers (\texttt{bar.sync} instructions), ensuring that the GEMMs (specifically, GEMM1 -- $\mathbf{P} \mathbf{V}$ from the current iteration, and GEMM0 -- $\mathbf{Q} \mathbf{K}^\top$ from the subsequent iteration) in warpgroup 1 are completed before those in warpgroup 2 begin. This strategy allows the softmax computation for warpgroup 1 to execute concurrently while warpgroup 2 carries out its GEMMs. Afterwards, the tasks switch: warpgroup 2 executes softmax, and warpgroup 1 proceeds with GEMMs—resulting in what is termed “pingpong” scheduling. This scheduling strategy is depicted in~\cref{fig:pingpong_scheduling}. Although the real implementation of pingpong scheduling does not match the idealized version shown in the figure, in practice it yields notable performance gains (for instance, throughput improvements from 570 TFLOPS to 620–640 TFLOPS for FP16 forward with a head dimension of 128 and sequence length of 8192).

\paragraph{Attention variants} In the case of multi-query attention~\citep{shazeer2019fast} and grouped query attention~\citep{ainslie2023gqa}, our method aligns with that utilized in \textsc{FlashAttention-2}, modifying tensor indexing so as to prevent redundancy of $\mathbf{K}$ and $\mathbf{V}$ within HBM.

\subsection{Intra-warpgroup overlapping GEMMs and softmax}

Within a single warpgroup, it is possible to concurrently execute certain softmax instructions alongside some GEMM operations. One approach to achieving this overlap is outlined below.

In the attention algorithm, operations within the inner loop (main loop) have sequential dependencies that impede parallelization within a single iteration. For example, (local) softmax (lines \ref{code-ws:softmax_start} to \ref{code-ws:softmax_end}) relies on the output $\mathbf{S}_i^{(j)}$ of the first GEMM, while the second GEMM takes its result $\widetilde{\mathbf{P}}_i^{(j)}$ as an operand. Indeed, the wait statements in lines \ref{alg:ws_only_gemm1} and \ref{alg:ws_only_gemm2} of \cref{alg:flash3_wgmma_ws_only} serialize the execution of softmax and GEMMs. However, we can break these dependencies by pipelining across iterations through additional buffers in registers. Pursuing this idea, we propose the following two-stage\footnote{Note that the number of stages of the overlapping scheme is bounded by, but need not equal, the number $s$ of stages in the circular SMEM buffer.} GEMM-softmax pipelining algorithm:

\begin{algorithm}[H]
    \caption{\small\label{alg:flash3_wgmma}\textsc{FlashAttention-3} consumer warpgroup forward pass}
    \begin{algorithmic}[1]
      \REQUIRE  Matrices $\mathbf{Q}_i \in \mathbb{R}^{B_r \times d}$ and $\mathbf{K}, \mathbf{V} \in \mathbb{R}^{N \times d}$ residing in HBM, key block size $B_c$ with $T_c = \lceil \frac{N}{B_c} \rceil$.
      \STATE Adjust the allocation of registers according to the number of consumer warps present.
      \STATE Initialize on-chip values: $\mathbf{O}_i = (0) \in \mathbb{R}^{B_r \times d}$, $\ell_i, m_i = (0), (-\infty) \in \mathbb{R}^{B_r}$.
      \STATE Await loading of $\mathbf{Q}_i$ and $\mathbf{K}_0$ into shared memory.
      \STATE Perform matrix multiplication $\mathbf{S}_{\mathrm{cur}} = \mathbf{Q}_i \mathbf{K}_0^T$ via WGMMA. Ensure completion of operation.
      \STATE Release buffer space allocated for the initial stage for $\mathbf{K}$.
      \STATE Derive $m_{i}$, $\tilde{\mathbf{P}}_{\mathrm{cur}}$, and $\ell_{i}$ from $\mathbf{S}_{\mathrm{cur}}$, then update $\mathbf{O}_i$ via rescaling.
      \FOR{$1 \le j < T_c - 1$}
        \STATE Wait until $\mathbf{K}_j$ has been transferred to shared memory.
        \label{code:mainloop_start}
        \STATE Apply WGMMA to compute $\mathbf{S}_{\mathrm{next}} = \mathbf{Q}_i \mathbf{K}_{j}^T$. Initiate and continue without awaiting completion. \label{code:first_wgmma}
        \STATE Await $\mathbf{V}_{j-1}$ availability in shared memory.
        \STATE Update $\mathbf{O}_{i}$ with the term $\tilde{\mathbf{P}}_{\mathrm{cur}} \mathbf{V}_{j-1}$ via WGMMA. Initiate the computation, proceed without blocking. \label{code:second_wgmma}
        \STATE Ensure WGMMA operation $\mathbf{Q}_i \mathbf{K}_{j}^T$ has finalized.
        \STATE Calculate $m_{i}$, $\tilde{\mathbf{P}}_{\mathrm{next}}$, and $\ell_{i}$ using $\mathbf{S}_{\mathrm{next}}$. \label{code:softmax}
        \STATE After confirming completion of $\tilde{\mathbf{P}}_{\mathrm{cur}} \mathbf{V}_{j-1}$, rescale $\mathbf{O}_i$.
        \STATE Free $(j\,\%\,s)$th and $(j-1\,\%\,s)$th stages of the buffer designated for $\mathbf{K}$ and $\mathbf{V}$, respectively.
        \STATE Move contents of $\mathbf{S}_{\mathrm{next}}$ into $\mathbf{S}_{\mathrm{cur}}$.
        \label{code:mainloop_end}
      \ENDFOR \STATE Await $\mathbf{V}_{T_c - 1}$ to be available in shared memory.
      \STATE Compute
      $\mathbf{O}_{i} = \mathbf{O}_{i} + \tilde{\mathbf{P}}_{\mathrm{last}} \mathbf{V}_{T_c - 1}$ using WGMMA. Launch and synchronize for completion.

\STATE Finalization: Adjust $\mathbf{O}_{i}$ proportionally to $m_{i}$. Derive $L_{i}$ using $m_{i}$ alongside $\ell_{i}$.
    Store $\mathbf{O}_{i}$ and $L_{i}$ in HBM as the $i$-th segment within $\mathbf{O}$ and $L$.
  \end{algorithmic}
\end{algorithm}

\cref{alg:flash3_wgmma} serves as an alternative implementation for the consumer pathway in \cref{alg:flash3_wgmma_ws_only}, thereby forming the entire \textsc{FlashAttention-3} workflow for FP16 calculations. In this context, WGMMA is used as shorthand for asynchronous GEMM operations. During the main loop (spanning lines \ref{code:mainloop_start} through \ref{code:mainloop_end}), the second invocation of WGMMA in iteration $j$ (line \ref{code:second_wgmma}) is scheduled concurrently with the softmax procedures of the subsequent iteration, $j+1$ (line \ref{code:softmax}).

Although the pipelined architecture depicted earlier is promising from a theoretical standpoint, several pragmatic considerations must be addressed:

\paragraph{Compiler reordering} The pseudocode above assumes a perfect execution sequence. However, in practice, the NVCC compiler frequently reorders instructions for efficiency reasons. This reordering can interrupt the intended interleaving of WGMMA and non-WGMMA operations, potentially reducing performance benefits or resulting in unforeseen execution patterns. Examination of the resulting SASS code indicates that the expected instruction overlap is preserved (see Section~\ref{sec:2-stage-sass}).

\paragraph{Register pressure} Achieving maximum performance hinges on minimizing register spilling. The two-stage pipeline, however, imposes greater register requirements, as additional temporary values and contextual information between stages must be maintained. In particular, a supplementary $\mathbf{S}_{\mathrm{next}}$ must reside in registers, causing an incremental register footprint of $B_r \times B_c \times \text{sizeof}(\text{float})$ per threadblock. This elevated demand for registers can interfere with block-size enlargement—another optimization tactic that also increases register consumption. Actual performance decisions should therefore be guided by profiling and empirical results.

\paragraph{3-stage pipelining} Building upon the two-stage method, a three-stage pipelining approach is also proposed, enabling further overlap between the second WGMMA and softmax computation. While this has the potential to enhance Tensor Core usage, it also necessitates greater register allocation due to the extra stage, thus complicating the balance between tile dimension and pipeline depth. Full algorithmic details and evaluation for the three-stage variant appear in ~\cref{sec:3-stage}.

\subsection{Low-precision with FP8}
\label{sec:algofp8}

\textbf{Efficiency: layout transformations.} Executing the forward pass of \textsc{FlashAttention-3} using FP8 precision introduces unique difficulties related to layout compatibility, which are not present when working with FP16.

To begin with, it is important to recognize that the input tensors $\mathbf{Q}$, $\mathbf{K}$, and $\mathbf{V}$ are commonly organized with contiguous elements along the head dimension. However, to comply with the k-major requirement imposed by FP8 WGMMA during the second GEMM, we must ensure that $\mathbf{V}$—specifically, its tiles stored in SMEM—are contiguous in the sequence length dimension instead. Given that TMA loads do not modify the axis of contiguity, we are faced with two alternatives: (1) perform a transpose of $\mathbf{V}$ in GMEM as an initial preprocessing step, or (2) transpose the tiles within the kernel itself after transferring them to SMEM. For approach (1), this could be accomplished by (1a) integrating the transpose operation with the epilogue of an earlier procedure, such as rotary embedding, or (1b) invoking an independent kernel dedicated to transpose preprocessing\footnote{Optimized transpose kernels are capable of achieving performance close to the device’s memory bandwidth~\citep{colfax_cutlass_transpose_2024}.} that switches the stride order between sequence length and head dimensions. Nevertheless, integrating strategy (1a) into a typical library poses significant challenges, while employing strategy (1b) incurs excessive overhead in memory-constrained contexts like inference.

For FP8 \textsc{FlashAttention-3}, we choose to implement option (2). During the in-kernel transpose operation, we utilize the LDSM (\verb|ldmatrix|) and STSM (\verb|stmatrix|) instructions, which facilitate a warp of threads to collaboratively transfer data between SMEM and RMEM, as well as vice versa, with a chunk size of 128 bytes.\footnote{While LDSM/STSM are documented in PTX as operating on $8 \times 8$ matrices composed of 16-bit elements \cite[\S 9.7.13.4.15-16]{ptx}, packing two 8-bit values enables their use for FP8 computations. Nonetheless, the transpose variants of these instructions are unable to individually separate the packed 8-bit values, which requires explicit register manipulation between LDSM and STSM calls to accomplish the tile-wise transpose; we omit a detailed exposition here.} These instructions are highly efficient in their register usage, permitting execution within the producer warpgroup, and they also support memory copy operations that include transposing layouts. Additionally, beginning with the second iteration, we can schedule the transpose of the upcoming $\mathbf{V}$ tile to execute concurrently with the two WGMMAs that operate on the previous $\mathbf{V}$ and the present $\mathbf{K}$ tile, thereby overlapping computation and data movement.

Secondly, it is notable that the FP32 accumulator associated with an FP8 WGMMA possesses a memory layout distinct from that of operand A when it is stored in registers, a contrast not observed in FP16. Illustrative excerpts of these two layouts are provided in \cref{fig:rmem_accum} and \cref{fig:rmem_operand}, showing the ordering of register-held entries per thread. Through the application of byte permute instructions, the accumulator resulting from the initial WGMMA can be reformatted to suit the requirements of the subsequent WGMMA, in accordance with the organization of the $\mathbf{V}$ tile yielded by the in-kernel transpose. Precisely, considering \cref{fig:rmem_accum}, the transformation imposes a new sequence:
$$\{ \verb|d0 d1 d4 d5 d2 d3 d6 d7| \},$$
and this register arrangement is repeated every 8 bytes. At the tile level, this technique effectively reorders the columns of the $\mathbf{P}$ tile (for instance, columns $0189$ are mapped to the initial four columns). To ensure WGMMA produces a valid output tile, the in-kernel transpose can be adjusted to output the $\mathbf{V}$ tile with a corresponding row permutation.\footnote{By leveraging the in-kernel transpose, we gain additional flexibility, thereby avoiding the necessity for shuffle instructions to reassign registers between threads, as previously discussed in~\citep{colfax_fp8_flashattention_2024}.}

\textbf{Accuracy: block quantization and incoherent processing.} The FP8 (e4m3) numerical format allocates just 3 bits for the mantissa and 4 bits to represent the exponent, which leads to greater rounding errors compared to FP16/BF16. Furthermore, in large-scale models, outlier values frequently occur~\citep{dettmers2208llm, sun2024massive} whose magnitudes surpass most of the distribution, introducing additional challenges for quantization. The conventional approach is to apply per-tensor scaling~\citep{micikevicius2022fp8}, assigning one scaling factor to each tensor (such as a single scalar for $\mathbf{Q}$, $\mathbf{K}$, and $\mathbf{V}$). To mitigate the accumulation of quantization errors in FP8-based attention, we incorporate two distinct strategies:

\begin{enumerate}

\item \textbf{Block quantization}: In this approach, a single scalar is maintained for each block, meaning that for each of $\mathbf{Q}$, $\mathbf{K}$, and $\mathbf{V}$, the tensor is divided into blocks of dimensions $B_r \times d$ or $B_c \times d$, and each block undergoes separate quantization. This block-wise quantization can be integrated seamlessly with pre-attention operations (such as rotary embedding) without incurring additional latency, since rotary embedding is limited by memory bandwidth. Given that the \textsc{FlashAttention-3} algorithm is inherently block-oriented, we are able to rescale every block of $\mathbf{S}$ in accordance with block quantization without incurring extra computational overhead.

\item \textbf{Incoherent processing}: To mitigate the influence of extreme values, we apply a random orthogonal transformation $\mathbf{M}$ to both $\mathbf{Q}$ and $\mathbf{K}$ prior to FP8 quantization. As $\mathbf{M}$ is orthogonal, it satisfies $\mathbf{M} \mathbf{M}^\top = I$, ensuring that $(\mathbf{Q} \mathbf{M}) (\mathbf{K} \mathbf{M})^\top = \mathbf{Q} \mathbf{K}^\top$. Thus, the attention mechanism remains unchanged. By randomly mixing the entries of $\mathbf{Q}$ and $\mathbf{K}$ through multiplication with $\mathbf{M}$, we effectively distribute outlier values, thereby minimizing quantization errors. In implementation, we adopt the methodology from \citet{chee2024quip} and \citet{tseng2024quip}, selecting $\mathbf{M}$ as a product of random diagonal matrices of $\pm 1$ and a Hadamard matrix. This composition allows multiplication in $O(d \log d)$ time rather than $O(d^2)$, and the process can be efficiently combined with rotary embedding without extra computational cost.

\end{enumerate}

Experimental evidence presented in \cref{sec:numerical_error} confirms that these methods lower numerical error by as much as 2.6$\times$.

\section{Empirical Validation}
\label{sec:exp}

To implement \textsc{FlashAttention-3} and assess its performance and precision, we leverage CUTLASS~\citep{Thakkar_CUTLASS_2023} primitives, including abstractions like WGMMA and TMA.

\begin{itemize}

\item \textbf{Benchmarking attention.} We evaluate the performance of \textsc{FlashAttention-3} in terms of runtime across a range of sequence lengths and compare it against several baselines: the standard PyTorch implementation, \textsc{FlashAttention-2}, a version of \textsc{FlashAttention-2} leveraging Triton (utilizing H100-specific instructions), and a vendor-provided, cuDNN-optimized \textsc{FlashAttention-2} for H100 GPUs. Our results show that \textsc{FlashAttention-3} achieves up to a 2.0$\times$ speedup over \textsc{FlashAttention-2}, and a 1.5$\times$ improvement over the Triton variant of \textsc{FlashAttention-2}. Furthermore, \textsc{FlashAttention-3} attains performance as high as 740 TFLOPs/s, corresponding to 75\% of the H100 GPU’s peak theoretical throughput.
\item \textbf{Ablation study.} We demonstrate, through ablation, that optimizations such as warp-specialization and the integration of GEMM-softmax pipelining are key contributors to the enhanced speed observed in \textsc{FlashAttention-3}.
\item \textbf{Accuracy of FP8 attention.} Through evaluation, we establish that implementing block quantization and incoherent processing effectively lowers the numerical error in FP8 \textsc{FlashAttention-3} by a factor of 2.6$\times$.
\end{itemize}

\subsection{Benchmarking Attention}
\label{subsec:benchmark_attn}

We evaluate the execution time of several attention mechanisms on an H100 80GB SXM5 GPU, considering configurations both with and without causal masking, and head dimensions set at either 64 or 128, using FP16 precision. The corresponding findings are presented in~\cref{fig:benchmark_attn_fwd} and~\cref{fig:benchmark_attn_bwd}. Our measurements indicate that \textsc{FlashAttention-3} achieves a speedup of roughly 1.5-2.0$\times$ over \textsc{FlashAttention-2} during the forward computation and approximately 1.5-1.75$\times$ during the backward computation. When benchmarked against a conventional attention implementation, \textsc{FlashAttention-3} demonstrates acceleration factors ranging from 3$\times$ up to 16$\times$. Remarkably, for sequence lengths at or exceeding 1{,}000, \textsc{FlashAttention-3} not only maintains these gains but also outperforms the proprietary cuDNN library, specifically engineered for H100 hardware.

\paragraph{Benchmark settings:} We vary the sequence length as 512, 1k, ..., 16k, and set batch size so that the total number of tokens is 16k. We set the hidden dimension to 2048, and head dimension to be either 64, 128, or 256 (i.e., 32 heads, 16 heads, or 8 heads). To calculate the FLOPs of the forward pass, we use:

\begin{equation*}
  4 \cdot \text{seqlen}^2 \cdot \text{head dimension} \cdot \text{number of heads}.
\end{equation*}

By implementing causal masking, the value is halved, reflecting that nearly half of the computations are performed. For estimating the FLOPs associated with the backward pass, we scale the FLOPs from the forward pass by a factor of 2.5. This multiplier arises because the forward computation involves two matrix multiplications, while the backward pass entails five matrix multiplications owing to recomputation.

Additionally, we assess the execution time of FP8 during the forward pass using comparable configurations. The outcomes for headdim 256 are presented in ~\cref{fig:benchmark_attn_fp8}, and a comprehensive overview of all results can be found in \cref{sec:benchmark_attn_fp8_full}.

\subsection{Ablation Study: 2-Stage Pipelining Experiments}
\label{sec:2-stage-pipelining-experiments}

We perform an ablation study on both the 2-stage WGMMA-softmax pipelining approach and warp-specialization for non-causal FP16 \textsc{FlashAttention-3}, maintaining constant parameter values $\{\text{batch}, \text{seqlen}, \text{nheads}, \text{hdim}\} = \{ 4, 8448, 16, 128\}$. The data presented in~\cref{table:ablation_pipelining} demonstrates that leveraging asynchrony through warp-specialization and enabling concurrency between the GEMM and softmax operations produces a marked performance improvement, boosting throughput from 570 to 661 TFLOPs.

\subsection{Numerical Error Validation}
\label{sec:numerical_error}

Given the ongoing attention to the numerical inaccuracies of \textsc{FlashAttention}~\citep{golden2024flash}, this study evaluates and contrasts \textsc{FlashAttention-2}, \textsc{FlashAttention-3}, and a conventional attention implementation with a reference FP64 version. To mimic atypical activations and features in LLMs~\citep{dettmers2208llm, sun2024massive}, we sample the components of $\mathbf{Q}, \mathbf{K}, \mathbf{V}$ from the following distribution:

\begin{equation*}
  \mathcal{N}(0, 1) + \mathcal{N}(0, 100) \cdot \mathrm{Bernoulli}(0.001).
\end{equation*}

Specifically, the entries are sampled from a normal distribution with mean zero and a standard deviation of 1, but for a randomly selected 0.1\% of entries, an additional independent term is included, drawn from a normal distribution with standard deviation 10. The root mean squared error (RMSE) is then reported in~\cref{table:numerical_error}. In FP16, both \textsc{FlashAttention-2} and \textsc{FlashAttention-3} yield an RMSE that is 1.7$\times$ lower than the standard approach because the intermediate softmax calculations are maintained in FP32 precision. For the baseline attention method in FP8, per-tensor scaling is used, the matrix multiplication accumulation is conducted in FP32, and softmax intermediate values are stored in FP16. Owing to block quantization and incoherent processing, \textsc{FlashAttention-3} implemented in FP8 demonstrates 2.6$\times$ higher accuracy than the baseline method.

\section{Dicussion, Limitations, Conclusion}
\label{sec:discussion}

Through the development of \textsc{FlashAttention-3}, we have shown that leveraging innovative programming methods and hardware advancements—specifically asynchrony and low-precision computation—can significantly enhance both the speed and precision of attention mechanisms. Our results indicate an acceleration of attention by 1.5–2.0$\times$ over \textsc{FlashAttention-2}, along with a 2.6$\times$ reduction in FP8 numerical errors when compared to traditional per-tensor quantization approaches. Several challenges remain to be addressed in upcoming work, including: improving the kernel for LLM inference scenarios, incorporating a persistent kernel approach within the FP8 implementation,\footnote{For benchmarking purposes, FP16 \textsc{FlashAttention-3} utilizes a persistent kernel and a load balancing mechanism, in contrast to FP8 \textsc{FlashAttention-3}, which lacks these features. This difference is a contributing factor to the lower performance of FP8 \textsc{FlashAttention-3} for shorter sequence lengths and with causal masking relative to FP8 cuDNN kernels.} and investigating the influence of reduced precision on attention during large-scale training runs. While our experiments targeted Hopper GPUs, we anticipate that the strategies employed here can be transferred to other acceleration hardware. We believe that advancements resulting in faster, more precise attention primitives may enable novel use cases for tasks requiring extended context.

\end{document}